% ======================================================
% SLIDE: Ý NGHĨA KHOA HỌC VÀ THỰC TIỄN
% ======================================================
\begin{frame}
    \frametitle{Ý nghĩa}

    \begin{BulletTight}
        \item \textbf{Ý nghĩa khoa học:}
        \begin{itemize}
            \setlength{\itemsep}{0.25em}
            \item Đóng góp tri thức mới về kiểm thử tự động dựa trên LLM và hệ đa tác tử.
            \item Đề xuất cơ chế phối hợp giữa các tác tử để tăng độ chính xác và hiệu quả kiểm thử.
            \item Làm rõ vai trò của HITL trong đảm bảo độ tin cậy và giảm ảo giác.
            \item Xây dựng framework có khả năng tái sử dụng cho nghiên cứu tiếp theo.
        \end{itemize}

        \vspace{0.6em}
        \item \textbf{Ý nghĩa thực tiễn:}
        \begin{itemize}
            \setlength{\itemsep}{0.25em}
            \item Phát triển công cụ kiểm thử tự động tích hợp vào DevOps/CI-CD.
            \item Giảm chi phí và công sức kiểm thử thủ công.
            \item Tăng độ bao phủ, chất lượng và tốc độ kiểm thử phần mềm.
            \item Góp phần thúc đẩy ứng dụng AI trong công nghiệp phần mềm.
        \end{itemize}


    \end{BulletTight}
\end{frame}